% Generated by artifacts/reference-comparison/analyze_reference.py.
\begin{table}[t]
\centering
\small
\setlength{\tabcolsep}{4pt}
\begin{tabular}{@{}lrrrrr@{}}
\toprule
Schedule & SPQR 32 & SPQR 64 & KEM 1 & KEM 32 & KEM 64\\
\midrule
\multicolumn{6}{l}{Reconstructed keys among 2,047 future sends}\\
Alternating & 172 & 88 & 4 & 66 & 130\\
Loss & 198 & 102 & 6 & 66 & 130\\
Reordering & 219 & 121 & 37 & 93 & 163\\
Duplication & 172 & 88 & 4 & 66 & 130\\
9:1 & 748 & 388 & 28 & 338 & 658\\
Offline then exchange & 424 & 342 & 260 & 322 & 386\\
One-way & 2047 & 2047 & 2047 & 2047 & 2047\\
\midrule
\multicolumn{6}{l}{Mean component bytes per send}\\
Alternating & 32.150 & 58.971 & 819.658 & 44.047 & 25.023\\
Loss & 31.957 & 58.336 & 819.564 & 44.047 & 23.520\\
Reordering & 31.876 & 61.795 & 817.999 & 38.927 & 23.058\\
Duplication & 32.150 & 58.971 & 819.658 & 44.047 & 25.023\\
9:1 & 23.508 & 40.531 & 819.381 & 21.031 & 13.417\\
Offline then exchange & 32.957 & 60.376 & 869.973 & 44.125 & 25.118\\
One-way & 38.875 & 70.875 & 1223.938 & 44.969 & 25.954\\
\bottomrule
\end{tabular}
\caption{The seven retained schedules at the same snapshot coordinate: endpoint $A$ immediately after its first send. SPQR 32 and 64 identify polynomial-chunk widths in bytes; KEM 1, 32 and 64 identify whole-KEM local-send intervals. Each column uses its recorded construction, authentication graph and retained-state interface. Bytes cover the component wire; reconstructed keys are exact predictor matches.}
\label{tab:reference-schedules}
\end{table}
